% =============================================================================
% Supervised Bayesian Matrix Factorization: Corrected Update Derivations
% Author:  Andrew Walther
% Date:    March 2026 (Revision of February 12, 2026 Document)
%
% Summary of revisions from the 2/12/26 document:
%  [R1] Section 1A -- Fixed dimension notation on Y, L, F, E
%  [R2] Eq. (22)   -- Fixed z^{-k}_i definition (f_{jk'} -> \beta_{k'})
%  [R3] Eqs.(32-35)-- Added explicit E_q[.] notation for second moments
%  [R4] Eqs.(47-48)-- Fixed EBNM subscript for F update (i -> j)
%  [R5] Eq. (62)   -- Added missing exponent: \beta_k -> \beta_k^2
%  [R6] Sec. 6     -- Added explicit EBNM(x_k, s_k) formulation for beta
%  [R7] Eq. (83)   -- Fixed sign error (+ -> -) in tau objective
%  [R8] Eq. (85)   -- Fixed subscript typo (l_{jk} -> l_{ik})
%
% This revision also expands every derivation section with the full
% intermediate algebraic steps that were omitted from the original document.
% =============================================================================

\documentclass[11pt]{article}
\usepackage[margin=1in]{geometry}
\usepackage{amsmath, amssymb, amsthm}
\usepackage{bm}
\usepackage{xcolor}
\usepackage{hyperref}
\usepackage{booktabs}
\usepackage{array}
\usepackage{cancel}
\usepackage{tcolorbox}
\tcbuselibrary{skins}

% ---- Custom Commands --------------------------------------------------------
% Posterior mean
\newcommand{\lbar}[2]{\bar{l}_{#1#2}}
\newcommand{\fbar}[2]{\bar{f}_{#1#2}}
\newcommand{\betabar}[1]{\bar{\beta}_{#1}}
% Second moment (E[x^2])
\newcommand{\lsec}[2]{\overline{l^2_{#1#2}}}
\newcommand{\fsec}[2]{\overline{f^2_{#1#2}}}
\newcommand{\betasec}[1]{\overline{\beta^2_{#1}}}
% Partial residuals
\newcommand{\Rk}[0]{R^{-k}}
\newcommand{\zk}[0]{z^{-k}}
% Revision callouts
\newcommand{\revised}[1]{{\color{red}\textbf{[REVISED]:} #1}}

% ---- Coloured boxes ---------------------------------------------------------
\newtcolorbox{correctionbox}{
  colback=red!5!white, colframe=red!50!black,
  title=Revision Note, fonttitle=\bfseries\small
}
\newtcolorbox{ebnmbox}{
  colback=blue!5!white, colframe=blue!40!black,
  title=EBNM Problem, fonttitle=\bfseries\small
}
\newtcolorbox{derivbox}{
  colback=gray!6!white, colframe=gray!50!black,
  title=Derivation Steps, fonttitle=\bfseries\small
}

% ---- Document ---------------------------------------------------------------
\title{\textbf{Supervised Bayesian Matrix Factorization}\\
       \large Corrected Update Derivations with Full Step-by-Step Algebra\\[4pt]
       \normalsize Revision of the February 12, 2026 Document}
\author{Andrew Walther}
\date{March 2026}

\begin{document}
\maketitle
\tableofcontents
\newpage

% =============================================================================
\section{Introduction}
% =============================================================================

This document presents the full derivations for the Supervised Bayesian Matrix
Factorization (Supervised MF) model. It covers the objective function
(ELBO), the second-order Taylor approximation of the Cox survival likelihood,
and the complete coordinate-ascent updates for the variational posteriors
$q(l_k)$, $q(f_k)$, $q(\beta_k)$, and the noise precision $\tau$.

Every algebraic step is shown explicitly.  Correction boxes flag the eight
errors found in the February~12, 2026 document (R1--R8).

\subsection{Update Procedure}

The general procedure for prior and posterior updates is:

\textbf{For} $b = 1, \ldots, B$:
\begin{enumerate}
  \item Compute the Cox Taylor working quantities $z_i$, $W_{ii}$.
  \item Iterate over $k = 1, \ldots, K$ for each of the following:
  \begin{enumerate}
    \item Update $q(l_k)$, $g(l_k)$ with all other parameters fixed
          (combines genomics and survival terms --- differs from standard EBMF).
    \item Update $q(f_k)$, $g(f_k)$ with all other parameters fixed
          (identical to the standard EBMF problem).
    \item Update $q(\beta_k)$, $g(\beta_k)$ with all other parameters fixed
          (survival term only --- differs from standard EBMF).
  \end{enumerate}
  \item Update $\tau_{ij}$ (identical to the standard EBMF problem).
\end{enumerate}

% =============================================================================
\section{Notation and Model Setup}
% =============================================================================

\subsection{Index Conventions}

\begin{center}
\begin{tabular}{cll}
\toprule
Symbol & Dimension & Meaning \\
\midrule
$n$ & scalar & number of samples (patients) \\
$p$ & scalar & number of features (genes / CpG sites) \\
$K$ & scalar & number of latent factors \\
$i$ & $1\ldots n$ & sample index \\
$j$ & $1\ldots p$ & feature index \\
$k, k'$ & $1\ldots K$ & factor index \\
$Y$ & $n\times p$ & observed genomics data matrix \\
$L$ & $n\times K$ & loading matrix; $l_{ik} =$ loading of sample $i$ on factor $k$ \\
$F$ & $p\times K$ & factor matrix; $f_{jk} =$ weight of feature $j$ on factor $k$ \\
$\bm{\beta}$ & $K\times 1$ & survival coefficient vector \\
$\tau$ & $p\times 1$ & column-specific noise precision (inverse noise variance) \\
$(t_i,\delta_i)$ & & survival time and event indicator for sample $i$ \\
\bottomrule
\end{tabular}
\end{center}

\subsection{Posterior Moment Notation}

Throughout, $q(\cdot)$ denotes the variational posterior and $g(\cdot)$ the
prior.  We carefully distinguish the \emph{posterior mean} from the
\emph{posterior second moment}:

\begin{align}
  \lbar{i}{k} &\;:=\; \mathbb{E}_q[l_{ik}]
                & &\text{(posterior mean)} \label{eq:lbar}\\
  \lsec{i}{k}  &\;:=\; \mathbb{E}_q[l_{ik}^2]
                 = \mathrm{Var}_q(l_{ik}) + \lbar{i}{k}^2
                & &\text{(posterior second moment)} \label{eq:lsec}
\end{align}

and analogously $\fbar{j}{k}$, $\fsec{j}{k}$, $\betabar{k}$,
$\betasec{k}$.  These two quantities are \textbf{not} equal whenever there is
posterior uncertainty.  This distinction is essential for the precision update
(Section~\ref{sec:tau}) and for the error-in-variables correction in the
$\beta$ update (Section~\ref{sec:beta}).

\subsection{Models}

\begin{correctionbox}
\textbf{[R1]} The 2/12/26 document wrote $Y = L_{n\times j}F'_{n\times k} + E_{k\times j}$,
which has inconsistent and impossible dimensions.  Corrected below.
\end{correctionbox}

\paragraph{Genomics model.}
\begin{equation}\label{eq:genomics-model}
  Y_{n\times p} \;=\; L_{n\times K}\,F^{\top}_{K\times p} \;+\; E_{n\times p},
  \qquad E_{ij} \overset{\text{iid}}{\sim} \mathcal{N}\!\left(0,\,\tau_j^{-1}\right)
\end{equation}
Element-wise: $Y_{ij} = \sum_{k=1}^K l_{ik} f_{jk} + \varepsilon_{ij}$.

\paragraph{Survival model (Cox proportional hazards).}
\begin{equation}\label{eq:survival-model}
  h(t_i) = h_0(t_i)\,\exp\!\left(\eta_i\right), \qquad
  \eta_i = \bm{l}_i^{\top}\bm{\beta} = \sum_{k=1}^K l_{ik}\beta_k
\end{equation}
where $h_0$ is an unspecified baseline hazard.

\paragraph{Hierarchical priors.}
\[
  l_{ik} \overset{\text{iid}}{\sim} g_l \in \mathcal{G}_l, \quad
  f_{jk} \overset{\text{iid}}{\sim} g_f \in \mathcal{G}_f, \quad
  \beta_k  \overset{\text{iid}}{\sim} g_\beta \in \mathcal{G}_\beta
\]
All three priors are learned from data via the Empirical Bayes Normal Means
(EBNM) framework.

% =============================================================================
\section{Objective Function (ELBO)}\label{sec:elbo}
% =============================================================================

The Evidence Lower Bound (ELBO) to be \emph{maximized} over posteriors
$(q_l, q_f, q_\beta)$, priors $(g_l, g_f, g_\beta)$, and precision $\tau$ is

\begin{align}\label{eq:elbo}
  \mathcal{F} \;=\;&
  \mathbb{E}_{q_l, q_f, q_\beta}\!\left[\log P(Y, t, \delta \mid L, F,\bm{\beta}, \tau)\right]
  \nonumber\\
  &+ \sum_k \mathbb{E}_{q_{l_k}}\!\left[\log\frac{g_{l_k}(l_k)}{q_{l_k}(l_k)}\right]
   + \sum_k \mathbb{E}_{q_{f_k}}\!\left[\log\frac{g_{f_k}(f_k)}{q_{f_k}(f_k)}\right]
   + \sum_k \mathbb{E}_{q_{\beta_k}}\!\left[\log\frac{g_{\beta_k}(\beta_k)}{q_{\beta_k}(\beta_k)}\right]
\end{align}

The joint log-likelihood decomposes as (conditional independence of $Y$ and
$(t,\delta)$ given $L$):
\begin{equation}\label{eq:lik-decomp}
  \log P(Y, t, \delta \mid L, F,\bm{\beta}, \tau)
  = \underbrace{\log P(Y \mid L, F, \tau)}_{\text{genomics}}
  + \underbrace{\log P(t, \delta \mid L, \bm{\beta})}_{\text{survival}}
\end{equation}

% =============================================================================
\section{Taylor Approximation of the Cox Partial Likelihood}\label{sec:taylor}
% =============================================================================

Because $\log P(t,\delta \mid L,\bm{\beta})$ is non-conjugate in $L$ and
$\bm{\beta}$, we use a second-order Taylor expansion to obtain a locally
Gaussian surrogate that can be handled by EBNM.

\subsection{The General Second-Order Taylor Expansion}

The second-order Taylor expansion of a function $f(x)$ around a point $x_0$
is:
\begin{equation}\label{eq:taylor-general}
  f(x) \;=\; f(x_0)
  \;+\; \underbrace{f'(x_0)}_{\text{gradient/score}}\,(x - x_0)
  \;+\; \frac{1}{2}\underbrace{f''(x_0)}_{\text{Hessian/info}}\,(x - x_0)^2
\end{equation}

\subsection{Applying the Expansion to the Cox Log-Likelihood}

Let $\ell(\bm{\eta}) = \log P(t, \delta \mid L, \bm{\beta})$ be the Cox partial
log-likelihood, where the linear predictor for sample $i$ is:
\begin{equation}\label{eq:linear-predictor}
  \eta_i \;=\; \bm{l}_i^{\top}\bm{\beta} \;=\; \sum_k l_{ik}\beta_k
\end{equation}

Let $\hat{\bm{\eta}}$ denote the current expansion point.  Define the score
vector $\mathbf{u}$ and the negative diagonal Hessian $W$ as:
\begin{align}
  u_i &= \frac{\partial \ell}{\partial \eta_i}\bigg|_{\hat{\bm{\eta}}}
        \quad\text{(Cox score for sample }i\text{)} \\
  W_{ii} &= -\frac{\partial^2 \ell}{\partial \eta_i^2}\bigg|_{\hat{\bm{\eta}}}
           \;\geq\; 0
           \quad\text{(negative diagonal Hessian, always non-negative)}
\end{align}

Applying the second-order expansion (Eq.~\eqref{eq:taylor-general}) to
$\ell(\bm{\eta})$ around $\hat{\bm{\eta}}$:

\begin{derivbox}
\begin{align}
\ell(\bm{\eta})
  &\;\approx\;
  \ell(\hat{\bm{\eta}})
  + (\bm{\eta} - \hat{\bm{\eta}})^{\top}\underbrace{\nabla\ell(\hat{\bm{\eta}})}_{\mathbf{u}}
  + \frac{1}{2}(\bm{\eta} - \hat{\bm{\eta}})^{\top}
    \underbrace{\nabla^2\ell(\hat{\bm{\eta}})}_{H\,=\,-W}
    (\bm{\eta} - \hat{\bm{\eta}}) \label{eq:taylor-vec}\\
  &\;=\;
  \ell(\hat{\bm{\eta}})
  + \sum_i (\eta_i - \hat{\eta}_i)\,u_i
  + \frac{1}{2}\sum_i (\eta_i - \hat{\eta}_i)^2\,H_{ii} \label{eq:taylor-scalar}
\end{align}

Since $H_{ii} = -W_{ii}$ (negative Hessian), substitute:

\begin{equation}
  = \ell(\hat{\bm{\eta}})
  + \sum_i (\eta_i - \hat{\eta}_i)\,u_i
  - \frac{1}{2}\sum_i (\eta_i - \hat{\eta}_i)^2\,W_{ii}
\end{equation}

Expand $(\eta_i - \hat{\eta}_i)^2 = \eta_i^2 - 2\eta_i\hat{\eta}_i + \hat{\eta}_i^2$:

\begin{equation}
  = \sum_i \left[
    u_i(\eta_i - \hat{\eta}_i)
    - \frac{1}{2}W_{ii}\left(\eta_i^2 - 2\eta_i\hat{\eta}_i + \hat{\eta}_i^2\right)
  \right] + C
\end{equation}

Distribute and retain only terms involving $\eta_i$ (constants w.r.t.\ $\eta_i$
absorbed into $C$):

\begin{equation}
  = \sum_i \left[
    u_i\eta_i
    - \frac{1}{2}W_{ii}\eta_i^2
    + W_{ii}\hat{\eta}_i\eta_i
  \right] + C \label{eq:taylor-distributed}
\end{equation}

Group into the quadratic form $\sum_i\!\left(a_i\eta_i^2 + b_i\eta_i\right)$,
where:
\begin{equation}
  a_i = -\tfrac{1}{2}W_{ii}, \qquad b_i = u_i + W_{ii}\hat{\eta}_i
\end{equation}

\textbf{Completing the square:} For any quadratic $a_i\eta_i^2 + b_i\eta_i$
with $a_i < 0$:
\begin{equation}
  \sum_i \left(a_i\eta_i^2 + b_i\eta_i\right)
  = \sum_i a_i\!\left(\eta_i + \frac{b_i}{2a_i}\right)^2 + C
  = \sum_i a_i\!\left(\eta_i - c_i\right)^2 + C
  \label{eq:complete-square-rule}
\end{equation}
where $c_i = -b_i/(2a_i)$.  Substituting $a_i = -\tfrac{1}{2}W_{ii}$ and
solving for $c_i$:
\begin{equation}
  c_i = -\frac{b_i}{2a_i}
      = -\frac{u_i + W_{ii}\hat{\eta}_i}{2 \cdot (-\tfrac{1}{2}W_{ii})}
      = \frac{u_i + W_{ii}\hat{\eta}_i}{W_{ii}}
      = \frac{u_i}{W_{ii}} + \hat{\eta}_i \label{eq:ci-solve}
\end{equation}

Apply the complete-the-square form back to Eq.~\eqref{eq:taylor-distributed}:
\begin{align}
  &= \sum_i \left(-\frac{1}{2}W_{ii}\right)\!\left(\eta_i - c_i\right)^2 + C
  \label{eq:taylor-ctsq}\\
  &= \sum_i \left(-\frac{1}{2}W_{ii}\right)\!
     \left(\eta_i^2 - 2\eta_i c_i + c_i^2\right) + C
  \label{eq:taylor-expand-back}
\end{align}

\end{derivbox}

\subsection{The Working Response}

Collecting results, the Cox log-likelihood is approximated as:

\begin{equation}\label{eq:taylor-quad}
  \ell(\bm{\eta}) \;\approx\;
  \sum_i \left[-\frac{1}{2}\,W_{ii}\,(\eta_i - z_i)^2\right] + C
\end{equation}

where the \textbf{working response} is:
\begin{equation}\label{eq:working-response}
  \boxed{z_i \;=\; \hat{\eta}_i + \frac{u_i}{W_{ii}}}
\end{equation}

This is a weighted Gaussian likelihood with response $z_i$, precision $W_{ii}$,
and mean $\eta_i = \sum_k l_{ik}\beta_k$.  The quadratic surrogate is a
proper likelihood for a linear model, enabling EBNM to be applied directly to
both $L$ and $\bm{\beta}$.

% =============================================================================
\section{Update for $q_{l_k}$ --- Patient Loadings}\label{sec:update-L}
% =============================================================================

\subsection{Setting Up the Objective}

To update $q(l_k)$ and $g(l_k)$, we extract from the full ELBO only the terms
that depend on $l_k$ (or $l_k^2$):

\begin{align}
  \mathcal{F}(q_l, g_l)
  &= \mathbb{E}_{q_l}\!\left[\log P(Y, t, \delta \mid L, F, \bm{\beta}, \tau)\right]
    + \sum_k \mathbb{E}_{q_{l_k}}\!\left[\log\frac{g_{l_k}(l_k)}{q_{l_k}(l_k)}\right]
  \label{eq:Fql-a}\\
  &= \underbrace{\log P(Y \mid L, F, \tau)}_{\text{Eq.\ 1: Genomics}}
    + \underbrace{\log P(t, \delta \mid L, \bm{\beta})}_{\text{Eq.\ 2: Survival}}
    + \underbrace{\sum_k \mathbb{E}_{q_{l_k}}\!\left[\log\frac{g_{l_k}(l_k)}{q_{l_k}(l_k)}\right]}_{\text{KL terms (kept)}}
  \label{eq:Fql-b}
\end{align}

We need to identify the coefficients $A_{ik}$ and $B_{ik}$ attached to the
terms $l_{ik}^2$ and $l_{ik}$ in Equations~1 and~2, then combine them.

\subsection{Defining Partial Residuals}

\paragraph{Genomics partial residual.}
The full residual for entry $(i,j)$ is:
\begin{align}
  R_{ij} &= Y_{ij} - \sum_k l_{ik}f_{jk} \label{eq:full-resid}\\
         &= Y_{ij} - \left[l_{ik}f_{jk} + \sum_{k'\neq k}l_{ik'}f_{jk'}\right]
         \label{eq:resid-split}
\end{align}

Rearranging to isolate the $k$-th factor:
\begin{equation}\label{eq:R-residual}
  R_{ij} = \underbrace{\left(Y_{ij} - \sum_{k'\neq k}\lbar{i}{k'}\fbar{j}{k'}\right)}_{R^{-k}_{ij}}
           - l_{ik}f_{jk}
\end{equation}

The \textbf{$k$-th genomics partial residual} $R^{-k}_{ij}$ is the data residual
after removing all factors \emph{except} $k$:
\begin{equation}
  R^{-k}_{ij} = Y_{ij} - \sum_{k'\neq k}\lbar{i}{k'}\fbar{j}{k'}
  \quad\Rightarrow\quad
  Y_{ij} \;\approx\; R^{-k}_{ij} + l_{ik}f_{jk}
\end{equation}

\paragraph{Survival partial working response.}

\begin{correctionbox}
\textbf{[R2]} The 2/12/26 document (Eq.\ 22) wrote
$z^{-k}_i = z_i - \sum_{k'\neq k} l_{ik'}f_{jk'}$,
incorrectly placing the factor weight $f_{jk'}$ (which carries a
\emph{feature} index $j$) in a \emph{per-sample} quantity.  The working
response $z_i$ is a survival quantity and does not involve $F$ at all --- it
involves $\bm{\beta}$.  The correct definition is:
\end{correctionbox}

\begin{equation}\label{eq:z-residual}
  z^{-k}_i = z_i - \sum_{k'\neq k}\lbar{i}{k'}\betabar{k'}
  \quad\Rightarrow\quad
  z_i \;\approx\; z^{-k}_i + l_{ik}\beta_k
\end{equation}

\subsection{Genomics Likelihood Contribution (Equation~1)}

The genomics log-likelihood for sample $i$, feature $j$ is:
\begin{equation}
  \log P(Y_{ij} \mid l_{ik}, f_{jk}, \tau_j)
  = -\frac{1}{2}\tau_j\left(Y_{ij} - \sum_k l_{ik}f_{jk}\right)^2 + C
\end{equation}

\begin{derivbox}
Substituting the partial residual decomposition $Y_{ij} = R^{-k}_{ij} + l_{ik}f_{jk}$:
\begin{align}
  \log P(Y \mid L, F, \tau)
  &= \sum_{ij}\left[-\frac{1}{2}\tau_j\left(R^{-k}_{ij} - l_{ik}f_{jk}\right)^2\right]
  \label{eq:gen-lik-1}
\end{align}

Expand the squared term $(R^{-k}_{ij} - l_{ik}f_{jk})^2$:
\begin{align}
  &= \sum_{ij}\left[-\frac{1}{2}\tau_j
     \left(\cancel{\left(R^{-k}_{ij}\right)^2}
           - 2\,R^{-k}_{ij}\,l_{ik}\,f_{jk}
           + l_{ik}^2\,f_{jk}^2
     \right)\right]
  \label{eq:gen-lik-expand}
\end{align}

The $(R^{-k}_{ij})^2$ term is a constant with respect to $l_{ik}$ and is
absorbed into $C$.  Retaining only $l_{ik}$-dependent terms and taking the
expectation $\mathbb{E}_{q_f}[\cdot]$ over the posterior of $F$:
\begin{equation}
  \longrightarrow \sum_{ij}\left[
    \tau_j\,R^{-k}_{ij}\,l_{ik}\,\underbrace{\mathbb{E}_{q_f}[f_{jk}]}_{\fbar{j}{k}}
    - \frac{1}{2}\tau_j\,l_{ik}^2\,\underbrace{\mathbb{E}_{q_f}[f_{jk}^2]}_{\fsec{j}{k}}
  \right]
  \label{eq:gen-lik-terms}
\end{equation}
\end{derivbox}

Collecting the coefficients of $l_{ik}$ (term $B$) and $l_{ik}^2$ (term $A$)
from the genomics likelihood:

\begin{correctionbox}
\textbf{[R3]} The 2/12/26 document wrote $f_{jk}^2$ and $\beta_k^2$ as bare
squares.  These should be posterior second moments under $q_f$ and $q_\beta$:
$\mathbb{E}_{q_f}[f_{jk}^2] = \fsec{j}{k}$ and
$\mathbb{E}_{q_\beta}[\beta_k^2] = \betasec{k}$.  The distinction matters
because $\fsec{j}{k} = \mathrm{Var}_{q_f}(f_{jk}) + \fbar{j}{k}^2 \geq \fbar{j}{k}^2$.
\end{correctionbox}

\begin{equation}\label{eq:genomics-AB}
  A^{\text{gen}}_{ik} = \sum_j \tau_j\,\fsec{j}{k},
  \qquad
  B^{\text{gen}}_{ik} = \sum_j \tau_j\,R^{-k}_{ij}\,\fbar{j}{k}
\end{equation}

\subsection{Survival Likelihood Contribution (Equation~2)}

From the Taylor approximation (Eq.~\eqref{eq:taylor-quad}), with
$\eta_i = \sum_k l_{ik}\beta_k$:

\begin{derivbox}
Substitute the partial working response decomposition
$z_i = z^{-k}_i + l_{ik}\beta_k$, so $\eta_i - z_i = l_{ik}\beta_k - z^{-k}_i$:
\begin{align}
  \log P(t,\delta \mid L,\bm{\beta})
  &\approx \sum_i\left[-\frac{1}{2}W_{ii}\left(z^{-k}_i - l_{ik}\beta_k\right)^2\right]
  \label{eq:surv-lik-1}
\end{align}

Expand the squared term $(z^{-k}_i - l_{ik}\beta_k)^2$:
\begin{align}
  &= \sum_i\left[-\frac{1}{2}W_{ii}
     \left(
       \cancel{\left(z^{-k}_i\right)^2}
       - 2\,z^{-k}_i\,l_{ik}\,\beta_k
       + l_{ik}^2\,\beta_k^2
     \right)\right]
  \label{eq:surv-lik-expand}
\end{align}

The $(z^{-k}_i)^2$ term is constant w.r.t.\ $l_{ik}$.  Taking the expectation
$\mathbb{E}_{q_\beta}[\cdot]$ over the posterior of $\bm{\beta}$:
\begin{equation}
  \longrightarrow \sum_i\left[
    W_{ii}\,z^{-k}_i\,l_{ik}\,\underbrace{\mathbb{E}_{q_\beta}[\beta_k]}_{\betabar{k}}
    - \frac{1}{2}W_{ii}\,l_{ik}^2\,\underbrace{\mathbb{E}_{q_\beta}[\beta_k^2]}_{\betasec{k}}
  \right]
  \label{eq:surv-lik-terms}
\end{equation}
\end{derivbox}

Collecting coefficients of $l_{ik}$ and $l_{ik}^2$ from the survival term:
\begin{equation}\label{eq:survival-AB}
  A^{\text{surv}}_{ik} = W_{ii}\,\betasec{k},
  \qquad
  B^{\text{surv}}_{ik} = W_{ii}\,z^{-k}_i\,\betabar{k}
\end{equation}

\subsection{Combined Coefficients and EBNM Problem}

The full objective for $q(l_k)$ takes the form:
\[
  \mathcal{F}(q_{l_k}, g_{l_k})
  = \sum_k\left[
      -\frac{1}{2}A_{ik}\,l_{ik}^2 + B_{ik}\,l_{ik}
      + \mathbb{E}_{q_{l_k}}\!\left[\log\frac{g_{l_k}(l_k)}{q_{l_k}(l_k)}\right]
    \right]
\]

where the combined precision (quadratic coefficient) and sufficient statistic
(linear coefficient) are the sums of the genomics and survival contributions:

\begin{align}
  \boxed{A_{ik} = \sum_j \tau_j\,\fsec{j}{k} + W_{ii}\,\betasec{k}}
  \label{eq:A-L}\\[4pt]
  \boxed{B_{ik} = \sum_j \tau_j\,R^{-k}_{ij}\,\fbar{j}{k} + W_{ii}\,z^{-k}_i\,\betabar{k}}
  \label{eq:B-L}
\end{align}

This quadratic form $-\tfrac{1}{2}A_{ik}\,l_{ik}^2 + B_{ik}\,l_{ik}$ is
equivalent to a Gaussian likelihood $l_{ik} \mid x_i, s_i^2 \sim \mathcal{N}(x_i, s_i^2)$
with pseudo-observation $x_i = B_{ik}/A_{ik}$ and pseudo-variance
$s_i^2 = 1/A_{ik}$:

\begin{ebnmbox}
For each sample $i$, the update for factor $k$ solves the EBNM problem:
\begin{align}
  x_i &= \frac{B_{ik}}{A_{ik}}
       = \frac{\displaystyle\sum_j \tau_j\,R^{-k}_{ij}\,\fbar{j}{k}
               + W_{ii}\,z^{-k}_i\,\betabar{k}}
              {\displaystyle\sum_j \tau_j\,\fsec{j}{k} + W_{ii}\,\betasec{k}}
  \label{eq:xi-L}\\[6pt]
  s_i &= \frac{1}{\sqrt{A_{ik}}} \label{eq:si-L}\\[6pt]
  &(\hat{g}_{l_k},\, \hat{q}_{l_k}) \;=\; \mathrm{EBNM}(\mathbf{x}, \mathbf{s})
\end{align}
The EBNM solver returns the fitted prior $\hat{g}_{l_k}$ and posterior moments
$\lbar{i}{k} = \mathbb{E}_{\hat{q}}[l_{ik}]$ and
$\lsec{i}{k} = \mathrm{Var}_{\hat{q}}(l_{ik}) + \lbar{i}{k}^2$.
\end{ebnmbox}

% =============================================================================
\section{Update for $q_{f_k}$ --- Biological Factors}\label{sec:update-F}
% =============================================================================

\subsection{Setting Up the Objective}

$F$ appears \emph{only} in the genomics term of the likelihood (not in the
survival term), so the update is:

\begin{equation}\label{eq:Fqf}
  \mathcal{F}(q_{f_k}, g_{f_k})
  \;\approx\;
  \mathbb{E}_{q_l,q_f}\!\left[\log P(Y \mid L, F, \tau)\right]
  + \mathbb{E}_{q_{f_k}}\!\left[\log\frac{g_{f_k}(f_k)}{q_{f_k}(f_k)}\right]
\end{equation}

The genomics data model is:
\begin{equation}\label{eq:genomics-data-model}
  Y_{ij} \;\sim\; \mathcal{N}\!\left(\sum_k l_{ik}f_{jk},\; \tau_j^{-1}\right)
\end{equation}

\subsection{Deriving Coefficients for $f_{jk}$}

\begin{derivbox}
Expand the log-likelihood:
\begin{equation}
  \mathbb{E}_{q}\!\left[\log P(Y \mid L, F, \tau)\right]
  = \sum_{ij}\,\mathbb{E}_{q}\!\left[
    -\frac{1}{2}\tau_j\left(Y_{ij} - \sum_k l_{ik}f_{jk}\right)^2
  \right] \label{eq:Fqf-lik}
\end{equation}

Use the partial residual $R^{-k}_{ij} = Y_{ij} - \sum_{k'\neq k}\lbar{i}{k'}\fbar{j}{k'}$
so that $Y_{ij} = R^{-k}_{ij} + l_{ik}f_{jk}$:
\begin{equation}
  = \sum_{ij}\,\mathbb{E}_{q}\!\left[
    -\frac{1}{2}\tau_j\left(R^{-k}_{ij} - l_{ik}f_{jk}\right)^2
  \right] \label{eq:Fqf-resid}
\end{equation}

Expand the square $(R^{-k}_{ij} - l_{ik}f_{jk})^2$:
\begin{equation}
  = \sum_{ij}\,\mathbb{E}_{q}\!\left[
    -\frac{1}{2}\tau_j\left(
      \cancel{\left(R^{-k}_{ij}\right)^2}
      - 2\,R^{-k}_{ij}\,l_{ik}\,f_{jk}
      + l_{ik}^2\,f_{jk}^2
    \right)
  \right] \label{eq:Fqf-expand}
\end{equation}

The $(R^{-k}_{ij})^2$ term is constant w.r.t.\ $f_{jk}$ and drops.  Take the
expectation $\mathbb{E}_{q_l}[\cdot]$ over $L$, treating $R^{-k}_{ij}$ and
$\tau_j$ as fixed:
\begin{equation}
  = \sum_{ij}\left[
    \tau_j\,R^{-k}_{ij}\,\underbrace{\mathbb{E}_{q_l}[l_{ik}]}_{\lbar{i}{k}}\,f_{jk}
    - \frac{1}{2}\tau_j\,\underbrace{\mathbb{E}_{q_l}[l_{ik}^2]}_{\lsec{i}{k}}\,f_{jk}^2
  \right] \label{eq:Fqf-terms}
\end{equation}

Now group over $i$ (for each fixed $j$) to read off the quadratic form in $f_{jk}$:
\begin{equation}
  \mathcal{F}(q_{f_k}) \;\approx\;
  \sum_j\left[
    -\frac{1}{2}f_{jk}^2\underbrace{\left(\sum_i \tau_j\,\lsec{i}{k}\right)}_{A_{jk}}
    + f_{jk}\underbrace{\left(\sum_i \tau_j\,R^{-k}_{ij}\,\lbar{i}{k}\right)}_{B_{jk}}
  \right]
  + \mathbb{E}_{q_{f_k}}\!\left[\log\frac{g_{f_k}(f_k)}{q_{f_k}(f_k)}\right]
  \label{eq:Fqf-final}
\end{equation}
\end{derivbox}

\subsection{Coefficients}

Reading off from Eq.~\eqref{eq:Fqf-final}:

\begin{equation}\label{eq:AB-F}
  A_{jk} = \sum_i \tau_j\,\lsec{i}{k},
  \qquad
  B_{jk} = \sum_i \tau_j\,R^{-k}_{ij}\,\lbar{i}{k}
\end{equation}

\subsection{EBNM Problem}

\begin{correctionbox}
\textbf{[R4]} Equations 47--48 of the 2/12/26 document subscripted the EBNM
inputs as $x_i, s^2_i$ (sample index).  Since we produce one pseudo-observation
\emph{per feature} when updating factor $k$ (each $j$ provides one equation),
the correct subscript is $j$ (feature index).
\end{correctionbox}

\begin{ebnmbox}
For each feature $j$, solve the EBNM problem:
\begin{align}
  x_j &= \frac{B_{jk}}{A_{jk}}
       = \frac{\displaystyle\sum_i \tau_j\,R^{-k}_{ij}\,\lbar{i}{k}}
              {\displaystyle\sum_i \tau_j\,\lsec{i}{k}}
  \label{eq:xj-F}\\[6pt]
  s_j &= \frac{1}{\sqrt{A_{jk}}} \label{eq:sj-F}\\[6pt]
  &(\hat{g}_{f_k},\, \hat{q}_{f_k}) \;=\; \mathrm{EBNM}(\mathbf{x}, \mathbf{s})
\end{align}
\end{ebnmbox}

% =============================================================================
\section{Update for $q_{\beta_k}$ --- Survival Coefficients}
\label{sec:beta}
% =============================================================================

\subsection{Setting Up the Objective}

$\bm{\beta}$ appears only in the survival term. Dropping the genomics
log-likelihood (which is constant w.r.t.\ $\bm{\beta}$), the objective is:

\begin{equation}\label{eq:Fqb-setup}
  \mathcal{F}(q_\beta, g_\beta)
  = \mathbb{E}_{q_l, q_\beta}\!\left[\log P(t,\delta \mid L,\bm{\beta})\right]
  + \sum_k \mathbb{E}_{q_{\beta_k}}\!\left[\log\frac{g_{\beta_k}(\beta_k)}{q_{\beta_k}(\beta_k)}\right]
\end{equation}

\subsection{Preliminary: The Prior $g(\beta)$}

The prior on $\bm{\beta}$ is a zero-mean normal:
\begin{equation}
  \beta_k \;\overset{\text{iid}}{\sim}\; \mathcal{N}(0,\sigma^2),
  \qquad
  g(\beta_k) = \frac{1}{\sqrt{2\pi\sigma^2}}\exp\!\left(-\frac{\beta_k^2}{2\sigma^2}\right)
\end{equation}

Taking the expectation of the log-prior:
\begin{derivbox}
\begin{align}
  \mathbb{E}_{q_\beta}\!\left[\log g(\bm{\beta})\right]
  &= \mathbb{E}_{q_\beta}\!\left[
     \log\prod_k \frac{1}{\sqrt{2\pi\sigma^2}}\exp\!\left(-\frac{\beta_k^2}{2\sigma^2}\right)
     \right]
  \label{eq:log-prior-1}\\
  &= \mathbb{E}_{q_\beta}\!\left[
     \sum_k \left(
       \underbrace{\log\frac{1}{\sqrt{2\pi\sigma^2}}}_{\text{constant, drops}}
       - \frac{\beta_k^2}{2\sigma^2}
     \right)
     \right]
  \label{eq:log-prior-2}\\
  &= -\frac{1}{2\sigma^2}\,\mathbb{E}_{q_\beta}\!\left[\sum_k \beta_k^2\right]
  \label{eq:log-prior-3}
\end{align}
\end{derivbox}

\subsection{Preliminary: The Variance-Mean Decomposition}

A key identity used throughout is:
\begin{equation}\label{eq:var-mean-decomp}
  \mathbb{E}_q[l_{ik}^2]
  = \mathrm{Var}_q(l_{ik}) + \left(\mathbb{E}_q[l_{ik}]\right)^2
  = \mathrm{Var}_q(l_{ik}) + \lbar{i}{k}^2
\end{equation}

which means $\lsec{i}{k} = \mathrm{Var}_q(l_{ik}) + \lbar{i}{k}^2 \geq \lbar{i}{k}^2$.
Using the posterior means as shorthand:
$\mathbb{E}_q[l_{ik}] = \lbar{i}{k}$ and
$\mathbb{E}_q[l_{ik}^2] = \lsec{i}{k}$.

\subsection{Expanding the Survival Likelihood Over $L$}

From the Taylor approximation with $\eta_i = \sum_k l_{ik}\beta_k$:

\begin{equation}
  \mathcal{F}(q_\beta, g_\beta)
  = \mathbb{E}_{q_l, q_\beta}\!\left[
    \sum_i -\frac{1}{2}W_{ii}\left(z_i - \sum_k l_{ik}\beta_k\right)^2
  \right] + \text{(prior and posterior terms)}
\end{equation}

\begin{derivbox}
Expand the square $\left(z_i - \sum_k l_{ik}\beta_k\right)^2$:
\begin{align}
  &= \mathbb{E}_{q_l, q_\beta}\!\left[\sum_i -\frac{1}{2}W_{ii}
     \left(z_i^2 - 2z_i\sum_k l_{ik}\beta_k
           + \left(\sum_k l_{ik}\beta_k\right)^2
     \right)
  \right] \label{eq:Fqb-expand1}
\end{align}

Expand the cross term. For any vector product:
\begin{equation}
  \left(\sum_k l_{ik}\beta_k\right)^2
  = \sum_k l_{ik}^2\beta_k^2 + 2\sum_{k < k'} l_{ik}\beta_k\,l_{ik'}\beta_{k'}
  \label{eq:cross-expand}
\end{equation}

\begin{correctionbox}
\textbf{[R5]} The 2/12/26 document (Eq.\ 62) wrote
$(\sum_k l_{ik}\beta_k)^2 = \sum_k l_{ik}^2\beta_k + \ldots$,
missing the exponent on $\beta_k$.  The correct expansion is
$\sum_k l_{ik}^2\beta_k^{\mathbf{2}} + 2\sum_{k<k'}\ldots$ as shown above.
\end{correctionbox}

Substituting back:
\begin{align}
  &= \mathbb{E}_{q_l, q_\beta}\!\left[\sum_i -\frac{1}{2}W_{ii}
     \left(z_i^2 - 2z_i\sum_k l_{ik}\beta_k
           + \sum_k l_{ik}^2\beta_k^2
           + 2\sum_{k<k'} l_{ik}\beta_k l_{ik'}\beta_{k'}
     \right)
  \right] \label{eq:Fqb-expand2}
\end{align}

Now apply $\mathbb{E}_{q_l}[\cdot]$ (treat $\bm{\beta}$ as random under $q_\beta$,
treat $L$ as random under $q_l$).  Under mean-field independence,
$l_{ik} \perp l_{ik'}$ for $k \neq k'$, so:
\begin{equation}
  \mathbb{E}_{q_l}[l_{ik}] = \lbar{i}{k},\quad
  \mathbb{E}_{q_l}[l_{ik}^2] = \lsec{i}{k},\quad
  \mathbb{E}_{q_l}[l_{ik}l_{ik'}] = \lbar{i}{k}\lbar{i}{k'} \text{ for } k\neq k'
\end{equation}

Applying $\mathbb{E}_{q_l}[\cdot]$, the $z_i^2$ term drops (constant w.r.t.\ $L$):
\begin{align}
  = \mathbb{E}_{q_\beta}\!\Biggl[\sum_i -\frac{1}{2}W_{ii}\Biggl(
    &-2z_i\sum_k \lbar{i}{k}\beta_k \nonumber\\
    &+\sum_k \lsec{i}{k}\beta_k^2
     + 2\sum_{k<k'}\lbar{i}{k}\beta_k\lbar{i}{k'}\beta_{k'}
  \Biggr)\Biggr] \label{eq:Fqb-Eql}
\end{align}

Recognise that the last two terms reconstitute the posterior-mean-based square:
\begin{align}
  \left(z_i - \sum_k \lbar{i}{k}\beta_k\right)^2
  &= z_i^2 - 2z_i\sum_k\lbar{i}{k}\beta_k
    + \sum_k\lbar{i}{k}^2\beta_k^2
    + 2\sum_{k<k'}\lbar{i}{k}\lbar{i}{k'}\beta_k\beta_{k'}
  \label{eq:verify-square}
\end{align}

comparing with Eq.~\eqref{eq:Fqb-Eql} we see an extra term:
\begin{equation}
  \sum_k\lsec{i}{k}\beta_k^2 - \sum_k\lbar{i}{k}^2\beta_k^2
  = \sum_k\mathrm{Var}_q(l_{ik})\,\beta_k^2
\end{equation}

Therefore (after dropping the constant $z_i^2$ term):
\begin{align}
  \mathbb{E}_{q_l,q_\beta}\!\left[\cdots\right]
  = \mathbb{E}_{q_\beta}\!\Biggl[
    &\sum_i -\frac{1}{2}W_{ii}\left(z_i - \sum_k\lbar{i}{k}\beta_k\right)^2 \nonumber\\
    &+ \sum_i -\frac{1}{2}W_{ii}\sum_k\mathrm{Var}_q(l_{ik})\,\beta_k^2
  \Biggr] \label{eq:Fqb-recollect}
\end{align}
\end{derivbox}

\subsection{Combining All Terms}

Substituting Eq.~\eqref{eq:log-prior-3} and Eq.~\eqref{eq:Fqb-recollect}
into the full objective Eq.~\eqref{eq:Fqb-setup}:

\begin{align}
  \mathcal{F}(q_\beta, g_\beta)
  &= \mathbb{E}_{q_\beta}\!\left[\sum_i -\frac{1}{2}W_{ii}
     \left(z_i - \sum_k\lbar{i}{k}\beta_k\right)^2\right]
  \label{eq:Fqb-term1}\\
  &\;\;- \mathbb{E}_{q_\beta}\!\left[
     \frac{1}{2}\sum_k\left(\sum_i W_{ii}\,\mathrm{Var}_q(l_{ik})\right)\beta_k^2
  \right]
  \label{eq:Fqb-term2}\\
  &\;\;- \frac{1}{2\sigma^2}\,\mathbb{E}_{q_\beta}\!\left[\sum_k\beta_k^2\right]
  \label{eq:Fqb-term3}\\
  &\;\;- \mathbb{E}_{q_\beta}\!\left[\log q(\bm{\beta})\right]
  \label{eq:Fqb-term4}
\end{align}

\begin{derivbox}
Combine terms \eqref{eq:Fqb-term2} and \eqref{eq:Fqb-term3} (both contain
$\beta_k^2$):
\begin{align}
  &-\frac{1}{2}\mathbb{E}_{q_\beta}\!\left[
    \sum_k\left(\sum_i W_{ii}\,\mathrm{Var}_q(l_{ik}) + \frac{1}{\sigma^2}\right)\beta_k^2
  \right] \label{eq:Fqb-combine}
\end{align}

The full objective is therefore:
\begin{align}
  \mathcal{F}(q_\beta, g_\beta)
  &= \mathbb{E}_{q_\beta}\!\left[
     \sum_i -\frac{1}{2}W_{ii}\left(z_i - \sum_k\lbar{i}{k}\beta_k\right)^2
  \right] \nonumber\\
  &\;\; -\frac{1}{2}\mathbb{E}_{q_\beta}\!\left[
     \sum_k\left(\sum_i W_{ii}\,\mathrm{Var}_q(l_{ik}) + \frac{1}{\sigma^2}\right)\beta_k^2
  \right]
  - \mathbb{E}_{q_\beta}\!\left[\log q(\bm{\beta})\right]
  \label{eq:Fqb-full}
\end{align}
\end{derivbox}

\subsection{Coordinate-Ascent Update for $\beta_k$}

Fix all $\beta_{k'}$ for $k' \neq k$.  Define the survival partial working
response (as in Eq.~\eqref{eq:z-residual}):
\begin{equation}
  z^{-k}_i = z_i - \sum_{k'\neq k}\lbar{i}{k'}\betabar{k'}
\end{equation}

\begin{derivbox}
The first term in Eq.~\eqref{eq:Fqb-full} contributes, for factor $k$:
\begin{align}
  &\sum_i -\frac{1}{2}W_{ii}\left(z^{-k}_i - \lbar{i}{k}\beta_k\right)^2 \nonumber\\
  &= \sum_i -\frac{1}{2}W_{ii}\left(
    \cancel{\left(z^{-k}_i\right)^2}
    - 2\,z^{-k}_i\,\lbar{i}{k}\,\beta_k
    + \lbar{i}{k}^2\,\beta_k^2
  \right) \nonumber\\
  &\longrightarrow
    \underbrace{\left(\sum_i W_{ii}\,z^{-k}_i\,\lbar{i}{k}\right)}_{B_k}\,\beta_k
    - \frac{1}{2}\underbrace{\left(\sum_i W_{ii}\,\lbar{i}{k}^2\right)}_{\text{part of }A_k}\,\beta_k^2
  \label{eq:Fqb-k-term1}
\end{align}

The variance-correction term (Eq.~\eqref{eq:Fqb-combine}) contributes, for factor $k$:
\begin{equation}
  -\frac{1}{2}\left(\sum_i W_{ii}\,\mathrm{Var}_q(l_{ik}) + \frac{1}{\sigma^2}\right)\beta_k^2
  \label{eq:Fqb-k-term2}
\end{equation}

Combining \eqref{eq:Fqb-k-term1} and \eqref{eq:Fqb-k-term2} and using
$\lsec{i}{k} = \mathrm{Var}_q(l_{ik}) + \lbar{i}{k}^2$:
\begin{align}
  &-\frac{1}{2}\left(\sum_i W_{ii}\,\lbar{i}{k}^2
                    + \sum_i W_{ii}\,\mathrm{Var}_q(l_{ik})
                    + \frac{1}{\sigma^2}\right)\beta_k^2
   + \left(\sum_i W_{ii}\,z^{-k}_i\,\lbar{i}{k}\right)\beta_k \nonumber\\
  &= -\frac{1}{2}\underbrace{\left(\sum_i W_{ii}\,\lsec{i}{k} + \frac{1}{\sigma^2}\right)}_{A_k}\beta_k^2
     + \underbrace{\left(\sum_i W_{ii}\,z^{-k}_i\,\lbar{i}{k}\right)}_{B_k}\beta_k
  \label{eq:Fqb-final-form}
\end{align}
\end{derivbox}

\begin{correctionbox}
\textbf{[R6]} The 2/12/26 document proposed a generic WLS solver for $\bm{\beta}$
without deriving the EBNM formulation.  The coordinate-ascent derivation above
(Eq.~\eqref{eq:Fqb-final-form}) shows that each $\beta_k$ update has the same
quadratic EBNM structure as the $L$ and $F$ updates.  Critically, $A_k$ uses
$\lsec{i}{k}$ (the full second moment including posterior variance) rather than
$\lbar{i}{k}^2$ (the square of the posterior mean).  This implements the
\emph{error-in-variables} correction: uncertainty in $L$ inflates the effective
noise for $\beta_k$, preventing over-fitting.
\end{correctionbox}

The combined precision and sufficient statistic are:

\begin{align}
  \boxed{A_k = \sum_i W_{ii}\,\lsec{i}{k} + \frac{1}{\sigma^2}}
  \label{eq:A-beta}\\[4pt]
  \boxed{B_k = \sum_i W_{ii}\,z^{-k}_i\,\lbar{i}{k}}
  \label{eq:B-beta}
\end{align}

\begin{ebnmbox}
For each factor $k$, solve the EBNM problem with a single pseudo-observation:
\begin{align}
  x_k &= \frac{B_k}{A_k}
       = \frac{\displaystyle\sum_i W_{ii}\,z^{-k}_i\,\lbar{i}{k}}
              {\displaystyle\sum_i W_{ii}\,\lsec{i}{k} + 1/\sigma^2}
  \label{eq:xk-beta}\\[6pt]
  s_k &= \frac{1}{\sqrt{A_k}} \label{eq:sk-beta}\\[6pt]
  &(\hat{g}_{\beta_k},\, \hat{q}_{\beta_k}) \;=\; \mathrm{EBNM}(x_k,\, s_k)
\end{align}
\textbf{Key remark.} In the code, the $1/\sigma^2$ prior term is subsumed into
the EBNM prior estimation --- the EBNM solver fits $g_{\beta_k}$ from data,
so no fixed $\sigma^2$ is required.  Using $\lsec{i}{k}$ rather than
$\lbar{i}{k}^2$ in $A_k$ encodes the error-in-variables correction.
\end{ebnmbox}

\subsection{Connection to Weighted Least Squares}

The first term of Eq.~\eqref{eq:Fqb-full} has the form of a weighted
regression:
\[
  z_i = \sum_k \lbar{i}{k}\,\beta_k + \epsilon_i,
  \qquad \epsilon_i \sim \mathcal{N}\!\left(0,\,W_{ii}^{-1}\right)
\]

The coordinate-ascent EBNM update derived above is the natural Bayesian
generalisation of this WLS step that: (i) accounts for posterior uncertainty in
$L$ via $\lsec{i}{k}$, and (ii) applies an EBNM shrinkage prior on $\beta_k$.

% =============================================================================
\section{Update for $\tau$ --- Noise Precision}\label{sec:tau}
% =============================================================================

\subsection{Objective}

The components of the ELBO that depend on $\tau$ are:

\begin{correctionbox}
\textbf{[R7]} Equation 83 of the 2/12/26 document wrote
$\tfrac{1}{2}\sum_{ij}[\log\tau_{ij} \mathbf{+} \tau_{ij}\bar{R}^2_{ij}]$.
The sign on the second term is wrong: maximizing
$\log\tau + \tau\bar{R}^2$ w.r.t.\ $\tau > 0$ gives no finite maximum
(the function is unbounded above as $\tau\to\infty$).
The correct expression has a \textbf{minus} sign.
\end{correctionbox}

\begin{equation}\label{eq:tau-obj}
  \mathcal{F}(\tau)
  = \mathbb{E}_{q_l,\, q_f}\!\left[
    \sum_{ij}\left(\frac{1}{2}\log\tau_j
    - \frac{1}{2}\tau_j\left(Y_{ij} - \sum_k l_{ik}f_{jk}\right)^{\!2}
    \right)\right] + C
  = \frac{1}{2}\sum_{ij}
    \left[\log\tau_j \;{\color{red}-}\; \tau_j\,\bar{R}^2_{ij}\right] + C
\end{equation}

\subsection{Expected Squared Residuals}

\begin{correctionbox}
\textbf{[R8]} Equation 85 of the 2/12/26 document wrote the last term as
$\sum_k \overline{l^2_{jk}}\,\overline{f^2_{jk}}$, using the subscript $jk$ on
the loading.  The loading carries a \emph{sample} index $i$, not $j$.
\end{correctionbox}

The expected squared residual $\bar{R}^2_{ij}$ requires care because both $L$
and $F$ are random under $q$.

\begin{derivbox}
Starting from the definition:
\begin{equation}
  \bar{R}^2_{ij} = \mathbb{E}_{q_l,q_f}\!\left[
    \left(Y_{ij} - \sum_k l_{ik}f_{jk}\right)^{\!2}
  \right] \label{eq:R2bar-def}
\end{equation}

Expand the square. Let $\mu_{ij} = \sum_k\lbar{i}{k}\fbar{j}{k}$ (the posterior
mean reconstruction). Write $l_{ik} = \lbar{i}{k} + \delta l_{ik}$ and
$f_{jk} = \fbar{j}{k} + \delta f_{jk}$ where $\mathbb{E}[\delta l_{ik}] = 0$
and $\mathbb{E}[\delta f_{jk}] = 0$.

Then:
\begin{align}
  Y_{ij} - \sum_k l_{ik}f_{jk}
  &= Y_{ij} - \mu_{ij}
    - \sum_k\left(\lbar{i}{k}\delta f_{jk} + \fbar{j}{k}\delta l_{ik}
                  + \delta l_{ik}\delta f_{jk}\right)
  \label{eq:resid-expand}
\end{align}

Taking the expectation of the square under mean-field ($l_{ik}\perp f_{jk'}$):
\begin{align}
  \bar{R}^2_{ij}
  &= \left(Y_{ij} - \mu_{ij}\right)^2
    + \sum_k\lbar{i}{k}^2\,\mathrm{Var}_q(f_{jk})
    + \sum_k\fbar{j}{k}^2\,\mathrm{Var}_q(l_{ik})
    + \sum_k\mathrm{Var}_q(l_{ik})\,\mathrm{Var}_q(f_{jk})
  \label{eq:R2bar-expand}\\
  &= \left(Y_{ij} - \sum_k\lbar{i}{k}\fbar{j}{k}\right)^{\!2}
    + \sum_k\!\left[\lsec{i}{k}\,\fsec{j}{k} - \lbar{i}{k}^2\,\fbar{j}{k}^2\right]
  \label{eq:R2bar-final}
\end{align}

The second step uses the identity
$\lsec{i}{k}\fsec{j}{k} - \lbar{i}{k}^2\fbar{j}{k}^2$
$= \lbar{i}{k}^2\mathrm{Var}_q(f_{jk}) + \fbar{j}{k}^2\mathrm{Var}_q(l_{ik})
   + \mathrm{Var}_q(l_{ik})\mathrm{Var}_q(f_{jk})$.
The first term in Eq.~\eqref{eq:R2bar-final} is the squared residual at
posterior means; the second is the \emph{variance inflation} correction.
\end{derivbox}

\begin{equation}\label{eq:expected-R2}
  \bar{R}^2_{ij}
  = \left(Y_{ij} - \sum_k \lbar{i}{k}\fbar{j}{k}\right)^{\!2}
    + \sum_k\!\left[\lsec{i\mathbf{k}}{}\;\fsec{j}{k}
               - \lbar{i}{k}^2\;\fbar{j}{k}^2\right]
\end{equation}

\subsection{Deriving the Precision Estimator}

\begin{derivbox}
Taking the derivative of $\mathcal{F}(\tau)$ in Eq.~\eqref{eq:tau-obj}
with respect to $\tau_j$ and setting to zero:
\begin{equation}
  \frac{\partial\mathcal{F}}{\partial\tau_j}
  = \frac{1}{2}\sum_i\left(\frac{1}{\tau_j} - \bar{R}^2_{ij}\right) = 0
  \label{eq:tau-deriv}
\end{equation}

Solving for $\tau_j$:
\begin{equation}
  \frac{n}{2\tau_j} = \frac{1}{2}\sum_i\bar{R}^2_{ij}
  \quad\Rightarrow\quad
  \hat{\tau}_j = \frac{n}{\displaystyle\sum_i\bar{R}^2_{ij}}
  \label{eq:tau-solve}
\end{equation}
\end{derivbox}

\subsection{Precision Estimators}

\begin{equation}\label{eq:tau-colspec}
  \text{Column-specific (}\tau_{ij} = \tau_j\text{):}\quad
  \hat{\tau}_j = \frac{N}{\displaystyle\sum_i \bar{R}^2_{ij}}
\end{equation}

\begin{equation}\label{eq:tau-constant}
  \text{Constant (}\tau_{ij} = \tau\text{):}\quad
  \hat{\tau} = \frac{NP}{\displaystyle\sum_{ij}\bar{R}^2_{ij}}
\end{equation}

In Eqs.~\eqref{eq:tau-colspec} and~\eqref{eq:tau-constant}:
\begin{itemize}
  \item $N$ = number of samples (rows of $Y$)
  \item $P$ = number of features (columns of $Y$)
  \item The column-specific form estimates precision from $N$ observations
        in column $j$
  \item The constant form averages across all $NP$ entries
\end{itemize}

% =============================================================================
\section{Complete Algorithm Summary}\label{sec:algorithm}
% =============================================================================

\begin{tcolorbox}[colback=gray!5, colframe=black!60,
                  title=Supervised Matrix Factorization --- CAVI Algorithm,
                  fonttitle=\bfseries]

\textbf{Inputs:} $Y_{n\times p}$, $(t_i, \delta_i)_{i=1}^n$, number of
factors $K$

\textbf{Initialize:} $\bar{L}, \bar{F}$ via rank-$K$ SVD of $Y$;
  $\bar{\bm{\beta}}$ via Cox on initial loadings;
  $\hat{\tau}$ via column-wise sample variance

\bigskip
\textbf{For} $b = 1, 2, \ldots$ \textbf{until convergence:}

\smallskip
\quad\textbf{1.\ Compute survival working quantities} (one Taylor expansion
per outer iteration):
\[
  \hat{\bm{\eta}} = \bar{L}\,\bar{\bm{\beta}}, \quad
  (u_i, W_{ii}) = \text{CoxTaylor}(\hat{\eta}_i, t_i, \delta_i), \quad
  z_i = \hat{\eta}_i + u_i/W_{ii}
\]

\smallskip
\quad\textbf{2.\ For} $k = 1, \ldots, K$:

\smallskip
\qquad\textbf{(a) Update} $q_{l_k}$, $g_{l_k}$ (patient loadings):
\begin{align*}
  R^{-k}_{ij} &= Y_{ij} - \textstyle\sum_{k'\neq k}\lbar{i}{k'}\fbar{j}{k'},
  \quad
  z^{-k}_i = z_i - \textstyle\sum_{k'\neq k}\lbar{i}{k'}\betabar{k'} \\
  A_{ik} &= \textstyle\sum_j \tau_j\,\fsec{j}{k} + W_{ii}\,\betasec{k},
  \quad
  B_{ik} = \textstyle\sum_j \tau_j\,R^{-k}_{ij}\,\fbar{j}{k}
           + W_{ii}\,z^{-k}_i\,\betabar{k} \\
  &(\lbar{i}{k},\,\lsec{i}{k})
    \;\leftarrow\; \mathrm{EBNM}(x_i = B_{ik}/A_{ik},\; s_i = 1/\sqrt{A_{ik}})
\end{align*}

\qquad\textbf{(b) Update} $q_{f_k}$, $g_{f_k}$ (biological factors):
\begin{align*}
  A_{jk} &= \textstyle\sum_i \tau_j\,\lsec{i}{k},
  \quad
  B_{jk} = \textstyle\sum_i \tau_j\,R^{-k}_{ij}\,\lbar{i}{k}\\
  &(\fbar{j}{k},\,\fsec{j}{k})
    \;\leftarrow\; \mathrm{EBNM}(x_j = B_{jk}/A_{jk},\; s_j = 1/\sqrt{A_{jk}})
\end{align*}

\qquad\textbf{(c) Update} $q_{\beta_k}$, $g_{\beta_k}$ (survival coefficients):
\begin{align*}
  z^{-k}_i &= z_i - \textstyle\sum_{k'\neq k}\lbar{i}{k'}\betabar{k'} \\
  A_k &= \textstyle\sum_i W_{ii}\,\lsec{i}{k},
  \quad
  B_k = \textstyle\sum_i W_{ii}\,z^{-k}_i\,\lbar{i}{k}\\
  &(\betabar{k},\,\betasec{k})
    \;\leftarrow\; \mathrm{EBNM}(x_k = B_k/A_k,\; s_k = 1/\sqrt{A_k})
\end{align*}

\smallskip
\quad\textbf{3.\ Update} $\hat{\tau}$ (column-specific noise precision):
\[
  \bar{R}^2_{ij} = \!\left(Y_{ij} - \textstyle\sum_k \lbar{i}{k}\fbar{j}{k}\right)^{\!2}
                 + \textstyle\sum_k\!\left(\lsec{i}{k}\,\fsec{j}{k}
                   - \lbar{i}{k}^2\fbar{j}{k}^2\right),
  \quad
  \hat{\tau}_j = \frac{N}{\textstyle\sum_i \bar{R}^2_{ij}}
\]

\smallskip
\quad\textbf{4.\ Check convergence:}
\[
  \Delta_L = \max_k\|\bar{l}_k^{\text{new}} - \bar{l}_k^{\text{old}}\|_\infty < \text{tol},
  \quad
  \Delta_\beta = \max_k|\betabar{k}^{\text{new}} - \betabar{k}^{\text{old}}| < \text{tol}
\]
\end{tcolorbox}

% =============================================================================
\section{Errata Table --- Summary of Revisions}\label{sec:errata}
% =============================================================================

\begin{center}
\begin{tabular}{clll}
\toprule
ID & Location & Error in 2/12/26 doc & Correction \\
\midrule
R1 & Sec.\ 1A, model     & Wrong dimension subscripts on $Y,L,F,E$
                          & Eq.~\eqref{eq:genomics-model} \\
R2 & Eq.\ 22             & $z^{-k}_i$ defined with $f_{jk'}$ instead of $\beta_{k'}$
                          & Eq.~\eqref{eq:z-residual} \\
R3 & Eqs.\ 32--35        & $f^2_{jk},\beta^2_k$ should be posterior second moments
                          & Eqs.~\eqref{eq:A-L}--\eqref{eq:B-L} \\
R4 & Eqs.\ 47--48        & EBNM inputs subscripted $i$; should be $j$
                          & Eqs.~\eqref{eq:xj-F}--\eqref{eq:sj-F} \\
R5 & Eq.\ 62             & $\beta_k$ missing square: should be $\beta_k^2$
                          & Eq.~\eqref{eq:cross-expand} \\
R6 & Sec.\ 6             & $\beta$ update as generic WLS; EBNM form missing
                          & Eqs.~\eqref{eq:A-beta}--\eqref{eq:B-beta} \\
R7 & Eq.\ 83             & Sign: $\log\tau + \tau\bar{R}^2$ should be $-\tau\bar{R}^2$
                          & Eq.~\eqref{eq:tau-obj} \\
R8 & Eq.\ 85             & $\overline{l^2_{jk}}$ wrong index $j$; should be $i$
                          & Eq.~\eqref{eq:expected-R2} \\
\bottomrule
\end{tabular}
\end{center}

% =============================================================================
\section{Notes on the Current R Implementation}\label{sec:code-notes}
% =============================================================================

\paragraph{Orthogonalization resets second moments.}
Every 10 iterations the code applies an SVD rotation to $\bar{L}$ for
identifiability, then resets \texttt{EL2 <- EL\^{}2} and
\texttt{EF2 <- EF\^{}2}.  This drops the posterior variance, setting
$\lsec{i}{k} \leftarrow \lbar{i}{k}^2$.  The correct second moment
is $\lbar{i}{k}^2 + \mathrm{Var}_q(l_{ik}) \geq \lbar{i}{k}^2$.  The variance
is recovered at the next EBNM call, so the algorithm remains consistent but
may transiently underestimate posterior uncertainty every 10 iterations.

\paragraph{Taylor expansion computed once per outer iteration.}
The working quantities $z_i$ and $W_{ii}$ are computed once per outer loop
before the inner $k$-loop, so they use the linear predictor
$\hat{\bm{\eta}}$ from the \emph{start} of each outer iteration.  Within the
inner loop, as $\bar{L}$ and $\bar{\bm{\beta}}$ are updated, the true linear
predictor changes but $z_i$ and $W_{ii}$ are not refreshed.  This is a
standard approximation in IRLS-within-VI algorithms and is acceptable; the
outer loop restores exactness at each iteration.

\paragraph{Single-observation EBNM for $\beta_k$.}
For each $k$, the $\beta_k$ EBNM call uses a single pseudo-observation
$x_k$.  The EBNM package requires at least one observation, but the
empirical prior $g_{\beta_k}$ is estimated from just this one point, making
the prior estimation highly variable.  For larger $K$, this could be
improved by pooling information across factors.

\end{document}
